\documentclass[10pt,a4paper]{amsart}
\usepackage[margin=19mm]{geometry}
\usepackage{amsmath,amssymb,mathtools}
\usepackage[scr=rsfs,cal=euler]{mathalfa}
\usepackage{xfrac}
\usepackage[unicode,hidelinks,bookmarks=false]{hyperref}
\newcommand{\ie}{i.e.\ }
\newcommand{\cf}{cf.\ }
\newcommand{\resp}{respectively\ }
\newcommand{\Subsection}{\subsection}
\providecommand{\nobreakdash}{\nobreak\hbox{-}}
\setlength{\emergencystretch}{2em}
\multlinegap0pt
\usepackage{amssymb}
\newtheorem{theorem}{Theorem}
\newtheorem{lemma}{Lemma}
\def \epsk {\epsilon_{n,k}}
\title{Induced/Incomparable versus Ramsey}
\author{Yair Caro, Zsolt Tuza, Christina Zarb}
\date{Source: arXiv:2605.22077v2 (CC BY 4.0)}
\begin{document}
\maketitle

\noindent\emph{Source and licence.} Induced/Incomparable versus Ramsey. Version arXiv:2605.22077v2 (CC BY 4.0), distributed by arXiv under the Creative Commons Attribution 4.0 International licence, http://creativecommons.org/licenses/by/4.0/, as recorded in the arXiv licence metadata for this exact version and shown on the abstract page https://arxiv.org/abs/2605.22077v2. Full licence notice: \texttt{licenses/arxiv-2605.22077-CC-BY-4.0.txt}.\par\smallskip

\noindent\textit{Editorial scope.} Counted unit: the article's theorem theorem (source0.tex lines 675--687). The article's complete original proof of it is delivered with it (source0.tex lines 689--773, one \texttt{proof} environment of 85 lines). No mathematics is written, repaired or re-derived: the statement and the proof are the article's own lines, in the article's own notation, and no retained line is substituted or re-flowed.  Retained uncounted as the source-local context the proof needs: lines 479--487; lines 644--648.  Nothing was omitted from any retained span, so the package carries no omission record.  No retained line contains a citation, so the wrapper carries no bibliography and no bibliography entry.  No retained line contains a figure environment, an image-inclusion command or drawing code, so no image is delivered and the package declares no asset dependency.  Editorial formatting only, in addition: an AMS \texttt{amsart} preamble that restates the article's own declarations and macros used by the retained lines, namely the environments \texttt{theorem}, \texttt{lemma} on separate counters, together with the standard packages those lines need.  The retained environments print in this document as Theorem 1, Lemma 1 in order of appearance, whereas the article prints the same items with the numbering of its own full document.  The complete native source of the article as distributed in the arXiv e-print for this exact version is bundled unchanged as \texttt{sources/201112/source0.tex}.
\begin{theorem}\label{treethree}
Let $H$ be a graph on $k\geq 3$ vertices with independence number $\alpha(H)$.  Then 
\begin{enumerate}
\item{$f(H) \geq R(H,K_{\alpha(H)}) - 1$.  }
\item{If $H$ is connected, then  $f(H)  \geq (\alpha(H) - 1)(k-1)$.}
\item{If $H$ is a tree $T$ on $k$ vertices, then $f(H) \geq (\alpha(T)-1)(k-1) \geq
 (k-1)(\left \lceil \frac{k}{2} \right \rceil -1)$}.
\end{enumerate}
\end{theorem}

\begin{lemma}
\label{l:linforest}
	Let $n\geq k\geq 5$ be integers.
	Then every $P_k$-free graph of order $n$ contains an induced linear forest of order $2\lceil \frac{n}{k-1} \rceil - \epsk$.
\end{lemma}	

	\begin{theorem}
	For all $k\ge 5$,
	
	$$
	  f(P_k) = \left\{
	    \begin{tabular}{lll}
	      $\dfrac{(k-1)^2}{2}$ && if\/ $k$ is \vspace{2ex} odd, \\
	      $\dfrac{(k-1)(k-2)}{2}+1$ & \qquad & if\/ $k$ is even.
	    \end{tabular}
	  \right.
	$$

	\end{theorem}

\begin{proof}
	The independence number of $P_k$ equals $\lceil k/2 \rceil$.
	So, as done in Theorem~\ref{treethree}, one can take a graph whose $\lceil k/2 \rceil-1$ (i.e., depending on the parity of $k$, either $(k-1)/2$ or $(k-2)/2$) components are $K_{k-1}$.
	This matches the claimed formula if $k$ is odd.
	For $k$ even, we can add a further isolated vertex.
	In all these graphs, every $k$-vertex induced subgraph contains a triangle (thus cannot be a subgraph of $P_k$), and is disconnected (hence $P_k$-free).
	This implies that $f(P_k)$ cannot be smaller than claimed.
	
	For an upper bound of the same value, let $G$ be a $P_k$-exact graph with $f(P_k)$ vertices.
	Should $G$ be $P_k$-free, we obtain by Lemma~\ref{l:linforest} that any graph with more vertices than the claimed formula contains an induced linear forest of order $k$, hence a proper subgraph of $P_k$.
	Indeed, for odd $k$,
	
	$$
	  2\left\lceil \frac{\frac{(k-1)^2}{2}+1}{k-1} \right\rceil = 2 \left\lceil \frac{k-1}{2} + \frac{1}{k-1} \right\rceil = k+1
	$$
	and for even $k$,
	
	$$
	  2\left\lceil \frac{\frac{(k-1)(k-2)}{2}+2}{k-1} \right\rceil = 2 \left\lceil \frac{k-2}{2} + \frac{2}{k-1} \right\rceil = k \,.
	$$	
	This contradicts our assumptions on $G$.
	So, $G$ contains one or more copies of $P_k$, all of which must be induced subgraphs.
	
	If $v_1\dots v_{k+1}$ is a path of order $k+1$ in $G$, then omitting any of its ends we obtain an induced $k$-path, but the omission of a vertex from the middle cannot yield a (disconnected) linear forest, hence the $(k+1)$-path actually induces a cycle.
	If there exists a further incident edge going out from the cycle, say $v_1w$, then for a similar reason, we obtain that $v_kw$ is also an edge.
	But then $v_{k+1}v_kwv_1v_2\dots v_{k-1}$ is a path with a chord $v_{k+1}v_1$, a contradiction.
	So $\{v_1,\dots,v_{k+1}\}$ induces a connected component.
	Then any further vertex with $v_1\dots v_{k-1}$ would induce disconnected a linear forest, therefore no such vertex can exist and $|G|=k+1$ would follow, far from being extremal if $k>4$.
	
	The final case, where the longest paths in $G$ have exactly $k$ vertices, is a little more complicated.
	Observe that
	
	$$
	  \frac{(k-1)(k-2)}{2}+1 < \frac{(k-1)^2}{2}
	$$
	holds for all $k>3$, so that the proof will be done if we show that $|G|$ is smaller than $\frac{1}{2}(k-1)(k-2)+1$ if $k\geq 5$.
	
	Let $P=v_1\dots v_k$ be a longest path; it is induced, and $P-v_i$ is a disconnected linear forest of order $k-1$ for every $2\leq i\leq k-1$.
	Thus, as $v_1$ and $v_k$ are of degree 1 by the non-extendability of $P$, every vertex $w$ not in $P$ has at least two neighors among the internal vertices of $P$.
	Choose a longest path $P'$ in $G-V(P)$, and let $w$ be an endpoint of $P$.
	Consider two neighbors of $w$ in $P$, say $v_i$ and $v_j$, $i<j$.
	We have $j\geq i+2$, otherwise $v_1\dots v_iwv_{i+1}\dots v_k$ would be a path longer than $P$.
	
	Observe that both $P'v_iv_{i+1}\dots v_k$ and $P'v_jv_{j-1}\dots v_1$ are non-induced paths, having at least one chord $v_jw$ and $v_iw$, respectively
	By assumption their order cannot exceed $k$, and since they are non-induced, they actually have order at most $k-1$.
	Consequently,
	
	\begin{eqnarray}
	 k-1 & \geq & |P'| + |P| - (i-1) \,, \nonumber \\
	 k-1 & \geq & |P'| + j \,. \nonumber
	\end{eqnarray}
	Substituting $|P|=k$ and summing up, we obtain
	
	$$
	  2k-2 \geq 2|P'|+k+(j-i)+1 \geq 2|P'|+k+3 \,.
	$$
	Hence,
	
	$$
	  |P'| \leq \left\lfloor \frac{k-5}{2} \right\rfloor .
	$$
	This means that the graph $G'=G-V(P)$ of order $|G|-k$ is $P_{\left\lfloor \frac{k-3}{2} \right\rfloor}$-free.
	(For $k=5$ we otained that $G\cong P_5$ must hold if $P_5\subseteq G$ and $G$ is a $P_5$-exact $P_6$-free graph!)
	On applying Lemma~\ref{l:linforest} it follows that the number of vertices in the largest induced linear forest $F$ in $G'$ is not smaller than
	
	$$
	  2 \left\lceil \frac{|G|-k}{\left\lfloor \frac{k-5}{2} \right\rfloor} \right\rceil - \epsilon_{|G|-k,\left\lfloor \frac{k-3}{2} \right\rfloor}
	  \geq
	  4 \cdot \frac{|G|-k}{k-5} - 1 \,.
	$$
	This $F$ can be extended to an induced disconnected linear forest in $G$ with at least two further vertices, namely $v_1$ and $v_k$.
	Therefore, $|F|\leq k-3$ because $G$ is $P_k$-exact.
	Consequently,
	
	$$
	  4 \cdot \frac{|G|-k}{k-5} - 1 \leq |F| \leq k-3 \,,
	$$
	$$
	  |G| \leq \frac{(k-2)(k-5)}{4} + k = \frac{k^2 - 3k + 10}{4}
	  <
	  \frac{(k-1)(k-2)}{2} + 1
	$$
	for all $k\geq 5$.
	This implies that a graph with maximum path length $k$ cannot be extremal for $f(P_k)$, hence completing the proof.	
\end{proof}

\end{document}
